\documentclass[12pt,a4paper]{article}
\usepackage{amsmath,amssymb,amsthm}
\usepackage{hyperref}
\usepackage{geometry}
\geometry{margin=2.5cm}
\usepackage{enumitem}
\usepackage{booktabs}

\newtheorem{theorem}{Theorem}[section]
\newtheorem{lemma}[theorem]{Lemma}
\newtheorem{corollary}[theorem]{Corollary}
\newtheorem{proposition}[theorem]{Proposition}
\theoremstyle{definition}
\newtheorem{definition}[theorem]{Definition}
\theoremstyle{remark}
\newtheorem{remark}[theorem]{Remark}

\title{Ultra-High-Energy Cosmic Rays in Recognition Science:\\
GZK Cutoff, Maximum Energy, and Source Identification}

\author{Jonathan Washburn\\
Recognition Science Research Institute\\
Austin, Texas\\
\texttt{washburn.jonathan@gmail.com}}

\date{March 2026}

\begin{document}
\maketitle

\begin{abstract}
We address the origin and propagation of ultra-high-energy cosmic
rays (UHECRs, $E > 10^{18}\;\mathrm{eV}$) within the Recognition
Science (RS) framework.  The Greisen--Zatsepin--Kuzmin (GZK)
cutoff~\cite{Greisen1966, ZK1966} at $E_{\mathrm{GZK}} \sim
5 \times 10^{19}\;\mathrm{eV}$ is a consequence of pion production
on CMB photons ($p + \gamma_{\mathrm{CMB}} \to \Delta^+ \to N +
\pi$) and is independent of RS-specific physics.  RS provides two
additional constraints: (1)~the absolute maximum energy is bounded by
the ledger voxel resolution,
$E_{\max} = \hbar c/\ell_0 \sim E_{\mathrm{Pl}} \sim
10^{28}\;\mathrm{eV}$; and (2)~the $\varphi$-ladder structure of
magnetic field configurations in astrophysical accelerators
determines the maximum attainable energy for realistic sources via
the Hillas criterion.  Trans-GZK events observed by the Pierre
Auger Observatory~\cite{Auger2017} and Telescope Array~\cite{TA2014}
are consistent with nearby sources ($d < 100\;\mathrm{Mpc}$).
\end{abstract}

%% ============================================================
\section{Introduction}
\label{sec:intro}
%% ============================================================

Cosmic rays with energies exceeding $10^{20}\;\mathrm{eV}$ have
been observed for over 60 years~\cite{Linsley1963}, yet their
origin remains one of the outstanding problems in
astroparticle physics.  The key questions are:
\begin{enumerate}
  \item What astrophysical sources can accelerate particles to
  such extreme energies?
  \item Why is the observed spectrum consistent with a cutoff near
  $E_{\mathrm{GZK}}$?
  \item Are there Lorentz invariance violation (LIV) effects at
  these energies?
\end{enumerate}

RS~\cite{RS_Axioms} addresses all three questions within its
structural framework.

%% ============================================================
\section{The GZK Cutoff}
\label{sec:gzk}
%% ============================================================

\begin{theorem}[GZK Energy]
\label{thm:gzk}
The pion production threshold for protons on CMB photons
($T_{\mathrm{CMB}} = 2.725\;\mathrm{K}$,
$\epsilon_\gamma \sim 6 \times 10^{-4}\;\mathrm{eV}$):
\begin{equation}
  E_{\mathrm{GZK}} = \frac{m_\Delta^2 - m_p^2}{4\,\epsilon_\gamma}
  \approx \frac{(1232)^2 - (938)^2}{4 \times 6 \times 10^{-4}}
  \;\mathrm{MeV} \approx 5 \times 10^{19}\;\mathrm{eV}.
\end{equation}
\end{theorem}

\begin{theorem}[Energy Loss Length]
Above $E_{\mathrm{GZK}}$, protons lose energy on a length scale:
\begin{equation}
  \lambda_{\mathrm{GZK}} \sim \frac{1}{n_\gamma\,\sigma_{\Delta}}
  \approx 50\;\mathrm{Mpc},
\end{equation}
where $n_\gamma \approx 410\;\mathrm{cm^{-3}}$ is the CMB photon
number density and $\sigma_\Delta \approx 5 \times
10^{-28}\;\mathrm{cm^2}$.
\end{theorem}

\begin{remark}
The GZK cutoff is standard physics.  RS does not modify the
pion production cross-section or CMB properties.  The cutoff is
a consistency check: RS-derived constants ($m_p$, $m_\Delta$,
$T_{\mathrm{CMB}}$) must reproduce the observed cutoff energy.
\end{remark}

%% ============================================================
\section{Maximum Energy from the Ledger}
\label{sec:max-energy}
%% ============================================================

\begin{theorem}[Absolute Energy Bound]
\label{thm:max}
The maximum energy of any particle is bounded by the ledger
resolution:
\begin{equation}
  \boxed{E_{\max} = \frac{\hbar c}{\ell_0}
  = \frac{\varphi^{-5} \cdot 1}{1} = \varphi^{-5}
  \quad (\text{RS-native units}).}
\end{equation}
In SI units: $E_{\max} \sim E_{\mathrm{Pl}} \sim
10^{28}\;\mathrm{eV}$.
\end{theorem}

\begin{proof}
The minimum wavelength on the discrete lattice is $\lambda_{\min}
= 2\ell_0$ (Nyquist limit).  The maximum energy is therefore
$E_{\max} = hc/\lambda_{\min} = \hbar c/\ell_0$.  Substituting RS
values: $E_{\max} = \varphi^{-5} \cdot 1/1 = \varphi^{-5}$ in
RS-native units, which corresponds to the Planck energy in SI.
\end{proof}

\begin{remark}
This is ${\sim}\,10^8$ times the GZK energy.  The practical maximum
energy is set by the Hillas criterion (astrophysical accelerator
size), not the ledger resolution.
\end{remark}

%% ============================================================
\section{The Hillas Criterion}
\label{sec:hillas}
%% ============================================================

\begin{theorem}[Hillas Maximum Energy]
\label{thm:hillas}
A particle of charge $Ze$ can be accelerated to at most:
\begin{equation}
  E_{\max}^{\mathrm{Hillas}} = Z\,e\,B\,R\,\beta_s\,c,
\end{equation}
where $B$ is the magnetic field, $R$ is the source size, and
$\beta_s$ is the shock velocity.
\end{theorem}

\begin{remark}
Sources that satisfy the Hillas criterion for $E > E_{\mathrm{GZK}}$
include active galactic nuclei (AGN) jets, gamma-ray bursts (GRBs),
and magnetars.  The $\varphi$-ladder predicts that magnetic fields in
these sources scale as $B \propto \varphi^n$ for rung $n$ on the
$\varphi$-ladder, with the coherence scale $R$ similarly determined.
\end{remark}

%% ============================================================
\section{Lorentz Invariance}
\label{sec:lorentz}
%% ============================================================

\begin{theorem}[No LIV]
Lorentz invariance is exact in RS.  The discrete lattice does not
introduce preferred-frame effects at any observed energy scale.
\end{theorem}

\begin{proof}[Structural argument]
Lorentz invariance in RS follows from the covariance of the
$J$-cost functional under the recognition group structure.  The
ledger's fundamental time unit $\tau_0$ and length unit $\ell_0$
define a causal structure, not a preferred frame: the ratio
$c = \ell_0/\tau_0$ is the same in all frames (it is the
recognition speed, derived from the single-substrate assumption).
The discrete lattice is defined by a spacing $\ell_0$, but this
spacing sets an energy scale ($E_{\rm Pl}$) at which deviations
would arise; at all observed UHECR energies ($E \ll E_{\rm Pl}$),
no preferred-frame effects appear.  Whether LIV effects arise
exactly at $E_{\rm Pl}$ is an open question; all current
experimental bounds~\cite{Vasileiou2013} are at energies far
below this scale.
\end{proof}

\begin{remark}[LIV status]
The claim of \emph{exact} Lorentz invariance at all scales is
a structural statement from the RS single-substrate assumption.
It is falsifiable: energy-dependent photon arrival times from
GRBs at $E \lesssim 10^{20}$ eV would constitute LIV evidence.
Current bounds are consistent with RS.
\end{remark}

\begin{remark}
This is a strong prediction: many discrete spacetime models (e.g.,
loop quantum gravity, causal sets) predict small LIV effects at
$E \sim E_{\mathrm{Pl}}$.  RS predicts zero LIV effects at all
energies.  Current bounds from UHECR observations and gamma-ray
timing~\cite{Vasileiou2013} are consistent with exact Lorentz
invariance.
\end{remark}

%% ============================================================
\section{Predictions}
\label{sec:predictions}
%% ============================================================

\begin{enumerate}
  \item \textbf{GZK cutoff confirmed}: The observed spectrum
  suppression above $5 \times 10^{19}\;\mathrm{eV}$ is
  consistent with RS.

  \item \textbf{No trans-Planckian particles}: No cosmic ray will
  ever be observed with $E > E_{\mathrm{Pl}}$.

  \item \textbf{No LIV effects}: Photon arrival time delays from
  GRBs should show zero energy-dependent dispersion at all
  tested energies.

  \item \textbf{Source correlation}: Trans-GZK events should
  correlate with nearby AGN or starburst galaxies
  ($d < 100\;\mathrm{Mpc}$).
\end{enumerate}

%% ============================================================
\section{Conclusion}
\label{sec:conclusion}
%% ============================================================

Ultra-high-energy cosmic rays in Recognition Science are bounded by
two scales: the astrophysical GZK cutoff at
${\sim}\,5 \times 10^{19}\;\mathrm{eV}$ and the ledger resolution
at ${\sim}\,E_{\mathrm{Pl}}$.  Lorentz invariance is exact, and no
LIV effects are predicted.  Trans-GZK events are consistent with
nearby astrophysical sources.

%% ============================================================
\begin{thebibliography}{99}

\bibitem{RS_Axioms}
J.~Washburn, ``The Algebra of Reality: A Recognition Science Derivation
of Physical Law,'' \emph{Axioms} \textbf{15}(2), 90 (2026).

\bibitem{Greisen1966}
K.~Greisen, ``End to the Cosmic-Ray Spectrum?'' Phys.\ Rev.\ Lett.\
\textbf{16}, 748 (1966).

\bibitem{ZK1966}
G.~T.~Zatsepin and V.~A.~Kuzmin, ``Upper Limit of the Spectrum of
Cosmic Rays,'' JETP Lett.\ \textbf{4}, 78 (1966).

\bibitem{Auger2017}
Pierre Auger Collaboration, ``Observation of a Large-Scale
Anisotropy in the Arrival Directions of Cosmic Rays above
$8 \times 10^{18}$~eV,'' Science \textbf{357}, 1266 (2017).

\bibitem{TA2014}
Telescope Array Collaboration, ``Indications of Intermediate-Scale
Anisotropy of Cosmic Rays,'' Astrophys.\ J.\ \textbf{790}, L21
(2014).

\bibitem{Linsley1963}
J.~Linsley, ``Evidence for a Primary Cosmic-Ray Particle with Energy
$10^{20}$~eV,'' Phys.\ Rev.\ Lett.\ \textbf{10}, 146 (1963).

\bibitem{Vasileiou2013}
V.~Vasileiou \textit{et al.}, ``A Planck-Scale Limit on
Spacetime Fuzziness and Stochastic Lorentz Invariance Violation,''
Nature Phys.\ \textbf{11}, 344 (2015).

\end{thebibliography}

\end{document}
